\documentclass[a4paper,12pt]{article}

\setlength{\textwidth}{15.0cm}
\setlength{\textheight}{24.0cm}
\setlength{\topmargin}{0cm}
\setlength{\headsep}{0cm}
\setlength{\headheight}{0cm}
\pagestyle{plain}

\usepackage[dvips]{epsfig}
\usepackage{tikz}

\usepackage{amsmath,amssymb,amsthm}
\newtheoremstyle{break}{4pt}{4pt}{}{}{\bfseries}{\vspace{2 pt}}{\newline}{}
\theoremstyle{break}
\newtheorem{theorem}{Theorem}[subsection]
\newtheorem{example}[theorem]{Example}
\newtheorem{definition}[theorem]{Definition}
\newtheorem{notation}[theorem]{Notation}
\newtheorem{lemma}[theorem]{Lemma}
\newtheorem{corollary}[theorem]{Corollary}
\newtheorem{technique}[theorem]{Technique}
\newtheorem{pythonn}[theorem]{Python Code}
\newtheorem{related}[theorem]{Related Work}

\usepackage{comment}
\usepackage{listings}
% Define custom colors for Python syntax highlighting
\definecolor{codebackground}{RGB}{242, 242, 242}
\definecolor{codecomment}{RGB}{106, 153, 85}
\definecolor{codekeyword}{RGB}{0, 0, 255}
\definecolor{codestring}{RGB}{170, 55, 241}

% Define custom Python style for syntax highlighting
\lstdefinestyle{custompython}{
    language=Python,
    backgroundcolor=\color{codebackground},
    basicstyle=\ttfamily\footnotesize,
    commentstyle=\color{codecomment}\itshape,
    keywordstyle=\color{codekeyword}\bfseries,
    stringstyle=\color{codestring},
    showstringspaces=false,
    breaklines=true,
    breakatwhitespace=true,
    tabsize=4,
    frame=tb,
    framesep=4pt,
    framerule=0.5pt,
    numbers=left,
    numbersep=10pt,
    numberstyle=\footnotesize\color{gray},
    xleftmargin=15pt,
    xrightmargin=5pt,
    aboveskip=5pt,
    belowskip=5pt,
    captionpos=b
}

\newcommand{\pythoncode}[1]{
    \hspace*{0cm}
    \vspace{-0.3cm}
    \lstinputlisting[style=custompython]{#1}}

% Define custom colors for Julia syntax highlighting
\definecolor{juliabackground}{RGB}{240, 230, 255} % Light purple background
\definecolor{juliacomment}{RGB}{128, 0, 128}     % Purple comments
\definecolor{juliakeyword}{RGB}{75, 0, 130}      % Indigo keywords
\definecolor{juliastring}{RGB}{147, 112, 219}    % Medium purple strings

% Define custom Julia style for syntax highlighting
\lstdefinelanguage{Julia}%
{
    morekeywords={
            abstract, break, case, catch, const, continue, do, else, elseif, end,
            export, false, for, function, global, if, immutable, import, importall,
            in, macro, module, otherwise, quote, return, switch, true, try, type,
            typealias, using, while
        },
    sensitive=true,
    alsoother={$},
    morecomment=[l]\#,
    morecomment=[n]{\#=}{=\#},
    morestring=[s]{"}{"},
    morestring=[s]{'}{'},
}[keywords,comments,strings]

\lstdefinestyle{customjulia}{
    language=Julia,
    backgroundcolor=\color{juliabackground},
    basicstyle=\ttfamily\footnotesize,
    commentstyle=\color{juliacomment}\itshape,
    keywordstyle=\color{juliakeyword}\bfseries,
    stringstyle=\color{juliastring},
    showstringspaces=false,
    breaklines=true,
    breakatwhitespace=true,
    tabsize=4,
    frame=tb,
    framesep=4pt,
    framerule=0.5pt,
    numbers=left,
    numbersep=10pt,
    numberstyle=\footnotesize\color{gray},
    xleftmargin=15pt,
    xrightmargin=5pt,
    aboveskip=5pt,
    belowskip=5pt,
    captionpos=b,
    columns=flexible,
    keepspaces=true
}

\newcommand{\juliacode}[1]{
    \hspace*{0cm}
    \vspace{-0.3cm}
    \lstinputlisting[style=customjulia]{#1}
}

\setlength{\parindent}{0pt}

\begin{document}

\begin{comment}
\title{Unbiased Monte Carlo for Recursive Integrals}
\author{Isidoor Pinillo Esquivel}
\maketitle
\end{comment}

\begin{titlepage}
    \begin{center}
        \resizebox{3cm}{!}{\includegraphics{./vert2_kl_01.eps}}
        \ \
        \ \\
        \ \\
        \ \\
        \ \\
        \ \\
        \ \\
        \ \\
        \ \\
        \ \\
        \ \\
        \ \\
        \Large{M{\sc asterproef scriptie}}
        \ \\
        \ \\
        \ \\
        \huge{\bf{\em Recursive Monte Carlo for linear ODEs}}
        \ \\
        \ \\
        \ \\
        % \ \\
        \ \\
        \ \\
        \normalsize
        Auteur: {\em Isidoor Pinillo Esquivel}\\
        \ \\
        \ \\
        Promotor: {\em Wim Vanroose}\\
        \ \\
        \ \\
        \ \\
        \ \\
        \ \\
        \ \\
        \ \\
        A{\sc cademiejaar 2022-2023}

    \end{center}
\end{titlepage}




\newpage
\tableofcontents
\newpage

\begin{abstract}
    We will write this at the end.
\end{abstract}

\section{Introduction}

\subsection{Introductory Example}
To get familiar with Monte Carlo for estimating recursive integrals
we demonstrate it on following problem:

\begin{equation} \label{ydy}
    y'=y, y(0)=1.
\end{equation}

By integrating both sides of (\ref{ydy}) following integral equation can be derived:

\begin{equation} \label{Integral ydy}
    y(t) = 1 + \int_{0}^{t} y(s) ds.
\end{equation}

Equation (\ref{Integral ydy}) is a recursive integral equation or to be more specific
a linear Volterra integral equation of the second type. By naively using Monte
Carlo on the recursive integral of equation (\ref{Integral ydy}) one derives following estimator:

\[
    Y(t) = 1 + t y(Ut)
    .\]

where $U=\text{Uniform}(0,1)$. If $y$ is well behaved then $E[Y(t)]=y(t)$ but we can't
calculate $Y(t)$ without accesses to $y(s),s<t$. Notice that we can replace $y$ by a
unbiased estimator of it without changing $E[Y(t)]=y(t)$ by the law of total expectance
($E[X] = E[E[X|Z]]$). By replacing $y$ by $Y$ itself we obtain a recursive expression for $Y$:

\begin{equation}\label{recursive RV}
    Y(t) = 1 + tY(Ut).
\end{equation}

Equation (\ref{recursive RV}) is a recursive random variable equation. If you would implement equation
(\ref{recursive RV}) with recursion it will run indefinitely. A biased way of around this is by approximating
$Y(t) \approx 1$ near $t = 0$. Later we discuss Russian roulette (ref{Russian  roulette}) which
can be used as an unbiased stopping mechanism. \\

\begin{python}[implementation of (\ref{recursive RV})] \label{python eps ydy}
    \pythoncode{python code/eps_ydy.py}
\end{python}

An issue with (\ref{python eps ydy}) is that the variance explodes when $t$ increases. Which we
later solve in the section on ODEs. Note that (\ref{python eps ydy}) keeps desirable properties
from unbiased Monte Carlo methods such as: being embarrassingly parallel,
robustness and simple error estimates.



\subsection{Contributions}
We write this at the end.

\subsection{Related Work}
work on
\begin{itemize}
    \item alternative methods for recursive integrals
    \item MC work on ODEs
    \item MC work on PDEs
    \item WoS
\end{itemize}
This is just to give a general overview we probably reference specific ideas when we first introduce them.

\section{Background}
\subsection{Modifying Monte Carlo}
Introduces Russian roulette, splitting, control variates, importance sampling and maybe quasi Monte Carlo with the
$y'=y$ example.

\subsection{Monte Carlo Trapezoidal Rule}
comparing normal vs Monte Carlo trapezoidal rule and highlighting the "half variance phenomenon".
+ maybe integrating polynomials for intuition

\subsection{Unbiased Non-Linearity}
\begin{example}[$y'=y^{2}$]
    see python note book
\end{example}

\begin{example}[$e^{E[X]}$]
    see python note book
\end{example}

\subsection{Recursion}

\begin{example}[coupled recursion]
    example with $y'=y$ (I need to redo this example)
\end{example}

\begin{example}[recursion in recursion]
    maybe induction in induction proof example
\end{example}

\begin{example}[tail recursion]
    discuss problems with implementing recursion and solutions.  \\
    inverse problem example
\end{example}

\subsection{Green Functions}
green function stuff that we will be needing, we aren't sure in how much detail we're going to go.

\begin{example}[numerical green functions]
    There will be probably some green functions that we need
    that don't have an analytic expression yet.
\end{example}

\section{$1$-Dimensional Recursive Integrals}

\subsection{Linear Recursive Integrals}
We have algo in mind for this case based on coupled recursion on disjunct sets.


\subsection{IVPs ODEs}
An IVP example probably using DRRMC maybe compare it to parareal. Maybe also non-linear algo

\subsection{BVPs ODEs}
A BVP example using yet another algo that hopefully has the half variance phenomenon.

\section{Higher Dimensional Recursive Integrals}
\subsection{Complicated Geometry}
\begin{example}[nasty $2$D integral]
    $2$D integral that is difficult because of its geometry
\end{example}

\subsection{Recursive Brownian Motion}
WoS like way to simulate Brownian motion which is related to the green function
of the heat equation

\begin{example}[recursive Brownian motion]
    see period5
\end{example}

\subsection{Heat Equation}
a geometric robust way to solve the heat equation and maybe a higher order method to solve
the heat equation

\subsection{Wave Equation}
probably won't get to it

\section{Appendix}
Derivation of the green functions and some expressions.

\end{document}
